\chapter{Results and Discussion}

\section{Performance Outcomes}
The culmination of the development, architectural design, and rigorous testing phases resulted in a highly capable and stable forensic platform. 

\subsection{UI Latency and Responsiveness}
One of the most significant achievements was the optimization of the frontend interface. Thanks to the architectural decisions regarding lazy loading, minimal API payloads, and asynchronous fetching, the UI achieved an exceptional performance metric: maintaining an average latency of \textbf{4.3 milliseconds} during data rendering. This ensures that security analysts receive near-instantaneous feedback when navigating complex forensic data, a crucial factor during time-sensitive incident response scenarios.

\subsection{Platform Stability and Cross-Platform Support}
The proactive resolution of Linux dependencies (\texttt{libpcre3}, \texttt{libpcap}, \texttt{python3-dev}) successfully enabled cross-platform functionality. Vigil-Core now operates seamlessly on both Windows and Linux, significantly broadening its deployability across various SOC infrastructures. Furthermore, the successful stress testing and implementation of chunked processing guarantee that the platform remains stable even when analyzing massive 8GB memory dumps.

\section{Analytical Accuracy}
The integration of specialized detection mechanisms yielded high accuracy in threat identification.
\begin{itemize}
    \item \textbf{YARA Validation:} The rigorous testing and refinement of YARA rules minimized false positives, ensuring that the alerts generated by the platform are highly relevant and actionable.
    \item \textbf{Fuzzy Hashing:} The implementation of \texttt{pydeep/ssdeep} successfully demonstrated the platform's ability to detect slightly modified malware variants, providing a layer of detection resilience that standard cryptographic hashing lacks.
\end{itemize}

\section{Discussion of Architectural Impact}
The decision to rely entirely on real, live data from the backend rather than using mock data for the UI was a defining architectural choice. While it required more intensive joint debugging to stabilize the API-UI connection, it ultimately resulted in a trustworthy, manager-ready dashboard. Complex forensic artifacts---such as entropy analysis, network protocol breakdowns, and extracted Indicators of Compromise (IoCs)---are now presented through clear, interactive visualizations, bridging the gap between raw data and strategic decision-making.
